% mage-auto-heading: the-capability
% mage-auto-heading: when-this-stack-earns-its-keep
% mage-auto-heading: the-composition
% mage-auto-heading: constituent-moves
% mage-auto-heading: why-these-travel-together
% Authoritative LaTeX manuscript; Markdown is a generated compatibility projection.

\section{The Capability}\label{heading:the-capability}

\textbf{Put relevant engineering knowledge and obligations into the actor\textquotesingle s context before the decision they must inform.}

\section{When This Stack Earns Its Keep}\label{heading:when-this-stack-earns-its-keep}

Reach for it when:

\begin{itemize}
\tightlist
\item
  \textbf{An actor must reason from engineering knowledge it does not carry} --- models, rules, prior decisions.
\item
  \textbf{The whole environment will not fit} in the actor\textquotesingle s working state, and dumping all of it into context merely relocates the navigation problem.
\item
  \textbf{Some obligations are critical and decidable}, and leaning on the actor to remember them is not good enough.
\end{itemize}

\section{The Composition}\label{heading:the-composition}

\begin{figure}[htbp]
\centering
\includegraphics[width=\linewidth]{book/latex/assets/context-delivery-stack-1b15a7023a.pdf}
\caption{The context-delivery composition. Engineering knowledge --- models, rules, decisions --- splits into a STANDING POLICY that is always seen and a TASK SLICE retrieved as needed; both feed the actor\textquotesingle s context, which drives the action. A dashed POINT-OF-ACTION attachment reasserts important obligations at the moment of action, and for critical decidable obligations becomes a deterministic gate rather than a reminder. Solid path: the essential composition. Dashed attachment: a useful enhancement, not required for the capability.}
\label{context-delivery-stack}
\end{figure}

\section{Constituent Moves}\label{heading:constituent-moves}

{\footnotesize
{\def\LTcaptype{none} % do not increment counter
\begin{longtable}[]{@{}
  >{\raggedright\arraybackslash}p{(\linewidth - 2\tabcolsep) * \real{0.5000}}
  >{\raggedright\arraybackslash}p{(\linewidth - 2\tabcolsep) * \real{0.5000}}@{}}
\toprule\noalign{}
\begin{minipage}[b]{\linewidth}\raggedright
Move
\end{minipage} & \begin{minipage}[b]{\linewidth}\raggedright
Role
\end{minipage} \\
\midrule\noalign{}
\endhead
\bottomrule\noalign{}
\endlastfoot
\textbf{STANDING} & Supply compact policy that applies broadly. \\
\textbf{SLICE} & Retrieve the task-relevant model and engineering context. \\
\textbf{DELIVER} & Put that context into the actor\textquotesingle s reasoning horizon. \\
\textbf{REINFORCE} & \emph{(strengthens)} Reassert important obligations at the point of action where useful. \\
\textbf{GATE} & For critical decidable obligations, do not rely on delivery at all --- enforce them. \\
\end{longtable}
}
}

\section{Why These Travel Together}\label{heading:why-these-travel-together}

An actor cannot reason through knowledge it never sees. But loading the entire engineered environment into every context only moves the navigation problem downstream: now the actor must find the relevant fact in a wall of them. Context delivery therefore combines compact standing policy with task-specific retrieval.

Point-of-action reminders can reinforce obligations that still require judgment. Critical decidable obligations should instead become deterministic controls: context is a reasoning aid, not enforcement. Boot files, dynamic snippets, hooks, and nudges are current implementations; the portable pattern is standing policy, task-specific retrieval, and deterministic enforcement where warranted.

\textbf{Mechanisms:} standing context · dynamic context injection · point-of-action guidance
